\section{已有 benchmark 证据}

\begin{frame}{已有结果证据的几条线}
\begin{itemize}
  \item 2026-04-22：初期 CPU benchmark 图表和摘要。
  \item 2026-04-28：12 charts presentation figures，区分 correctness / efficiency / memory。
  \item 2026-05-11：symbolic vs numeric efficiency evaluation。
  \item 2026-05-11/12/14：thread scaling 与 P/E-core profile。
  \item \code{results/cross-platform-v1/}：跨平台 schema 和 core-profile comparison。
\end{itemize}
\source{results/; source\_index.md}
\end{frame}
\note{
长期 deck 不重新发明数据，而是把结果包串起来。每次新增图表都应该记录来源，避免未来不记得某个数字来自哪次运行。
}

\begin{frame}{symbolic reuse vs direct no-symbolic：主结论}
\scriptsize
\begin{tabular}{rrrr}
\toprule
组装次数 & symbolic reuse 摊销总耗时 ms & direct no-symbolic 总耗时 ms & 收益 \\
\midrule
1  & 3818.786 & 6276.199 & \metric{1.644x} \\
3  & 1633.165 & 5822.369 & \metric{3.565x} \\
10 & 901.325  & 5696.788 & \metric{6.320x} \\
30 & 686.745  & 5552.854 & \metric{8.086x} \\
\bottomrule
\end{tabular}
\vspace{0.7em}

\takeaway{单次组装也受益；多次组装更能体现 symbolic result reuse 的摊销优势。}
\source{results/2026-05-11-symbolic-numeric/symbolic\_numeric\_eval\_report.md}
\end{frame}
\note{
这张表是 mentor 关心的核心结果。注意平台是 Apple M4 Max，它回答的是 symbolic reuse 是否有价值，不应该写成 Intel 绝对性能。主线比较是 \code{symbolic\_reuse\_serial} vs \code{direct\_no\_symbolic\_serial}。
}

\begin{frame}{2026-05-16：parallel symbolic vs direct/no-symbolic}
\scriptsize
\begin{tabular}{lrrrr}
\toprule
模式 & 线程 & symbolic total ms & numeric/direct ms & total ms \\
\midrule
\code{parallel\_symbolic\_reuse + cpu\_atomic} & 1  & 2829.454 & 544.796 & 3374.250 \\
\code{direct\_no\_symbolic\_parallel}          & 1  & 0        & 6571.854 & 6571.854 \\
\code{parallel\_symbolic\_reuse + cpu\_atomic} & 14 & 554.432  & 107.991 & 662.423 \\
\code{direct\_no\_symbolic\_parallel}          & 14 & 0        & 1669.417 & 1669.417 \\
\bottomrule
\end{tabular}
\vspace{0.7em}

\takeaway{同为 14 线程，parallel symbolic reuse 仍优于 direct/no-symbolic parallel；direct 路径主要输在 bucket/merge 与 sort/reduce。}
\source{results/2026-05-16-mentor-action-items/windhub\_parallel\_symbolic\_direct.md}
\end{frame}
\note{
这页是 mentor 2026-05-14 action items 的新证据。它不再只比较串行 symbolic/direct，而是在 WindHub 上跑了 1..14 物理核的 parallel symbolic 和 direct/no-symbolic parallel。14 线程时，parallel symbolic reuse 总计约 662 ms，direct/no-symbolic parallel 总计约 1669 ms。
}

\begin{frame}{为什么多次组装收益更明显}
\begin{itemize}
  \item symbolic 阶段成本主要来自 CSR pattern 和 scatter plan。
  \item 网格拓扑和 DOF 编号不变时，这部分可以只做一次。
  \item numeric 阶段随材料状态、迭代状态、时间步重复执行。
  \item 组装次数越多，symbolic 成本摊到每次上的比例越低。
\end{itemize}
\takeaway{多次组装场景下，二阶段组装不是微优化，而是改变成本结构。}
\end{frame}
\note{
可以用生活类比：先做路线规划，再重复开车。路线规划一次很贵，但如果同一条路线要走很多次，规划成本就被摊薄。这里路线就是 CSR/scatter，开车就是 numeric assembly。
}

\begin{frame}{控制实验：每次重建 symbolic 结构}
\scriptsize
\begin{tabular}{rrrr}
\toprule
组装次数 & symbolic 构建次数 & 摊销总耗时 ms & 相对 direct 收益 \\
\midrule
1  & 1  & 3813.269 & 1.646x \\
3  & 3  & 3735.218 & 1.559x \\
10 & 10 & 3743.918 & 1.522x \\
30 & 30 & 3394.580 & 1.636x \\
\bottomrule
\end{tabular}
\vspace{0.7em}

\takeaway{\code{symbolic\_rebuild\_serial} 说明：主收益来自符号结果复用，但 CSR/scatter 路径本身也比直接排序归并更合理。}
\source{results/2026-05-11-symbolic-numeric/symbolic\_numeric\_eval\_report.md}
\end{frame}
\note{
这不是推荐使用方式。它每轮都重建 CSR 和 scatter plan，用来隔离变量。如果有人误解“二阶段组装就是每次都拆成两步”，这张控制实验可以说明真正关键的是复用。
}

\begin{frame}{WindHub Tet4 physical correctness evidence}
\centering
\resultfig{../../results/2026-04-28-12charts-repeat3-threads1to14/presentation_charts_12_v2/04_correctness_heatmap_04_windhub_physics_tet4.png}
\source{results/2026-04-28-12charts-repeat3-threads1to14/presentation\_charts\_12\_v2; historical field \code{physics\_tet4}}
\end{frame}
\note{
长期 deck 中直接引用 results 图表，不复制到 assets。这张图用于提醒：算法比较必须先看 correctness。rel_l2 等错误指标不通过时，速度图没有意义。图名里的 physics_tet4 是历史结果包字段；按当前代码口径应理解为 WindHub Tet4/C3D4 物理刚度路径。
}

\begin{frame}{WindHub Tet4 physical efficiency evidence}
\centering
\resultfig{../../results/2026-04-28-12charts-repeat3-threads1to14/presentation_charts_12_v2/04_efficiency_grouped_bars_04_windhub_physics_tet4.png}
\source{results/2026-04-28-12charts-repeat3-threads1to14/presentation\_charts\_12\_v2; historical field \code{physics\_tet4}}
\end{frame}
\note{
效率图用于比较不同算法随线程变化的资源利用情况。长期 deck 中应保留这类图，但解释时不能只看“哪个柱子最高”，还要结合正确性和内存。
}

\begin{frame}{WindHub Tet4 physical memory evidence}
\centering
\resultfig{../../results/2026-04-28-12charts-repeat3-threads1to14/presentation_charts_12_v2/04_memory_heatmap_04_windhub_physics_tet4.png}
\source{results/2026-04-28-12charts-repeat3-threads1to14/presentation\_charts\_12\_v2; historical field \code{physics\_tet4}}
\end{frame}
\note{
内存图尤其重要，因为 private CSR、COO sort-reduce 这类算法可能用内存换同步成本。工程可用性不能只由时间决定。
}

\begin{frame}{WindHub sparse pattern evidence}
\centering
\resultfig{../../results/2026-05-16-mentor-action-items/sparse_pattern/windhub_physics_tet4_spy_python.png}
\source{results/2026-05-16-mentor-action-items/sparse\_pattern; generated from exported row,col pattern}
\end{frame}
\note{
这是 mentor 要求的 sparse pattern 证据。图不是附件示意图，而是从项目导出的 serial/parallel CSR row,col pattern 生成。为了让 685152 维、27502200 个非零位置可以渲染，图像层使用 occupancy raster；原始 CSV 和 MatrixMarket 文件仍保留完整 row,col pattern。
}

\begin{frame}{Sparse pattern: original order vs RCM order}
\begin{columns}[T,onlytextwidth]
\begin{column}{0.49\textwidth}
\centering
\resultfig[0.56]{../2026-05-22-weekly-meeting-beamer/assets/windhub_physics_tet4_visual_spy_original_raster.png}
\end{column}
\begin{column}{0.49\textwidth}
\centering
\resultfig[0.56]{../2026-05-22-weekly-meeting-beamer/assets/windhub_physics_tet4_visual_spy_rcm_raster.png}
\end{column}
\end{columns}
\takeaway{WindHub 原始 \code{.inp} 编号下呈现块状聚簇；同一 CSR pattern 经过 RCM 的 \code{K(p,p)} 可视化后呈现带状稀疏。}
\source{reports/2026-05-22-weekly-meeting-beamer/assets/windhub\_physics\_tet4\_visual\_metadata.json}
\end{frame}
\note{
这页吸收 2026-05-22 周会后的答疑需求。mentor 图里的 B(p,p) 是重排序结果；当前项目新增同口径 K(p,p) 可视化。这里的 RCM 只用于解释和显示，不改变 benchmark 的组装流程。
}

\begin{frame}{Exact local sparse window}
\begin{columns}[T,onlytextwidth]
\begin{column}{0.49\textwidth}
\centering
\resultfig[0.56]{../2026-05-22-weekly-meeting-beamer/assets/windhub_physics_tet4_visual_exact_window_serial.png}
\end{column}
\begin{column}{0.49\textwidth}
\centering
\resultfig[0.56]{../2026-05-22-weekly-meeting-beamer/assets/windhub_physics_tet4_visual_exact_window_auto_serial.png}
\end{column}
\end{columns}
\takeaway{完整矩阵不能 1:1 画成图片；局部窗口可以不做压缩，直接展示真实 \code{row,col} 非零坐标。}
\source{reports/2026-05-22-weekly-meeting-beamer/assets/windhub\_physics\_tet4\_visual\_exact\_window\_*.csv}
\end{frame}
\note{
这页记录“为什么不压缩”和“怎样不压缩”的边界。全局图必须使用 occupancy raster，局部窗口则保留真实坐标。自动窗口选择对角附近非零最密集的 4096 x 4096 区域。
}

\begin{frame}{lock guard vs atomic：同步原语对照}
\scriptsize
\begin{tabular}{lrrr}
\toprule
算法 & 最好线程 & 最好 assembly ms & extra memory bytes \\
\midrule
\code{cpu\_atomic} & 14 & 104.437 & 0 \\
\code{cpu\_private\_csr} & 8 & 105.041 & 1760140800 \\
\code{cpu\_lock\_guard} & 10 & 365.510 & 1760140800 \\
\bottomrule
\end{tabular}
\vspace{0.7em}

\takeaway{\code{cpu\_lock\_guard} 是有效 C++ lock baseline，但在当前 WindHub/M4 Max 上慢于 atomic。}
\source{results/2026-05-16-mentor-action-items/windhub\_lock\_vs\_atomic.md}
\end{frame}
\note{
这页是 action item 5 的证据。\code{lock\_guard} baseline correctness 通过，但加锁/解锁开销明显高于 OpenMP atomic update。extra memory 与 \code{private\_csr} 8 线程同为 1760140800 bytes 是因为前者按 nnz 个 mutex，后者按 8 份 nnz double，数值上恰好相同。
}

\begin{frame}{Intel full-host thread scaling}
\centering
\resultfig{../../results/2026-05-11-thread-scaling-linux-intel/figures/thread_scaling_bound_dashboard.png}
\source{results/2026-05-11-thread-scaling-linux-intel/figures}
\end{frame}
\note{
Intel Core Ultra 7 265KF 是后续面向主要用户群体的优先平台。full-host 图回答的是在整机可用核心上，不同算法随线程数如何扩展。
}

\begin{frame}{Intel P-core-only 与 E-core-only}
\centering
\resultfig{../../results/cross-platform-v1/figures/core_profile_speedup_comparison_intel_u7_265kf.png}
\source{results/cross-platform-v1/figures; results/2026-05-12-thread-scaling-linux-intel-hybrid-core-supplement.md}
\end{frame}
\note{
Intel 这组 P/E-core profile 通过 Linux taskset 限定 CPU 集合，属于比较硬的 affinity restriction。解释时不能说成单纯“P-core 就一定更快”或“E-core 更慢”，因为资源数量和线程范围也不同。
}

\begin{frame}{Apple M4 Max QoS-biased profile}
\centering
\resultfig{../../results/cross-platform-v1/figures/core_profile_speedup_comparison_apple_m4_max.png}
\source{results/cross-platform-v1/figures; results/2026-05-14-thread-scaling-macos-m4max-qos-supplement.md}
\end{frame}
\note{
Apple 的 performance/efficiency profile 是 QoS-biased sensitivity run，不是 Linux taskset 那种硬绑核。长期 deck 必须保留这个解释，避免把 Apple 和 Intel 的 profile 名字当成同一种实验。
}

\begin{frame}{benchmark figures 与 presentation charts 的区别}
\begin{tabularx}{\textwidth}{>{\bfseries}lXX}
\toprule
类型 & 目标 & 约束 \\
\midrule
benchmark/report figures & 保存真实结果和可追溯性 & 不为了好看改变数据口径 \\
presentation charts & 帮助听众快速理解结论 & 可以重新排版，但必须标注来源和边界 \\
长期 beamer figures & 个人复习和进度追踪 & 优先可追溯，必要时吸收 presentation 图 \\
\bottomrule
\end{tabularx}
\takeaway{图表可以美化，benchmark truth 不能被美化改写。}
\end{frame}
\note{
这是近期图表讨论形成的稳定原则。长期 deck 可以直接引用 presentation_charts_12_v2，因为这些图来自已有结果。但如果未来 presentation 版为了叙事重新排版，source_index 必须记录原始结果来源。
}
